\documentclass[11pt,a4paper]{article}

\usepackage[utf8]{inputenc}
\usepackage[T1]{fontenc}
\usepackage[brazilian]{babel}
\usepackage{amsmath,amssymb,amsthm,mathtools}
\usepackage{booktabs}
\usepackage{array}
\usepackage{longtable}
\usepackage{enumitem}
\usepackage{hyperref}
\usepackage{geometry}
\usepackage{xcolor}
\usepackage{listings}
\usepackage{microtype}

\geometry{margin=2.5cm}

\newtheorem{theorem}{Teorema}
\newtheorem{proposition}[theorem]{Proposição}
\newtheorem{lemma}[theorem]{Lema}
\newtheorem{corollary}[theorem]{Corolário}
\theoremstyle{definition}
\newtheorem{definition}[theorem]{Definição}
\newtheorem{remark}[theorem]{Observação}
\newtheorem{example}[theorem]{Exemplo}
\theoremstyle{remark}
\newtheorem{claim}[theorem]{Conjectura}

\newcommand{\Q}{\mathbb{Q}}
\newcommand{\R}{\mathbb{R}}
\newcommand{\N}{\mathbb{N}}
\newcommand{\C}{\mathbb{C}}
\newcommand{\Z}{\mathbb{Z}}
\newcommand{\lII}{\ell^2}
\newcommand{\supp}{\operatorname{supp}}
\newcommand{\op}{\mathrm{op}}
\newcommand{\Lin}{\mathcal{L}}
\newcommand{\D}{\mathcal{D}}
\newcommand{\E}{\mathcal{E}}
\newcommand{\qform}{Q}
\newcommand{\RH}{\mathrm{RH}}

\title{Como chegar à positividade geral de Weil\\[2pt]
\large Roteiro técnico a partir das auditorias ES8-P14/ES8-P16}
\author{Notas técnicas}
\date{04/11/2026}

\begin{document}
\maketitle

\tableofcontents

\section{O que precisa ser provado (e não apenas verificado)}

Chegar à positividade geral de Weil --- equivalentemente, a uma prova de
\RH{} --- exige uma mudança de escala: o que o ES8-P16 audita são
\textbf{truncamentos racionais finitos}; o critério de Weil é uma afirmação
sobre \textbf{todas as funções de teste admissíveis} (ou, equivalentemente,
sobre uma forma quadrática fechada num espaço de Hilbert infinito).

O critério de Weil, na forma quadrática, diz que
\[
\RH \quad\Longleftrightarrow\quad
Q_W(f)\ge 0
\ \ \text{para toda } f \text{ admissível (tipicamente } f=g*\tilde g\text{)}.
\]

Aqui \(Q_W\) é o funcional definido pela fórmula explícita: lado dos zeros
\(\sum_\rho \hat f(\rho)\) versus lado geométrico (primos \(+\) arquimediano
\(+\) polos). Os truncamentos \(Q^{(N)}\) do ES8-P16 são apenas restrições de
\(Q_W\) a subespaços de dimensão \(4,6,8\). Positividade em dimensão finita
\textbf{não implica} positividade em dimensão infinita --- e o próprio ES8-P16
formaliza por que: dois prolongamentos infinitos distintos coincidem no
prefixo \(8\times 8\).

Portanto, o primeiro passo lógico não é ``aumentar \(N\)'', mas
\textbf{definir a forma infinita e provar que ela é fechada e limitada
inferiormente}.

\section{Rota canônica: forma de Dirichlet e operador auto-adjunto}

A literatura recente mostra que a forma de Weil, numa janela multiplicativa
\(I_\lambda=[\lambda^{-1},\lambda]\), pode ser reescrita como uma
\textbf{forma de Dirichlet fechada}:
\[
\E_\lambda \;=\; \E_{\text{arq}} + \E_{\text{primos}},
\]
onde a parte arquimediana é uma energia de salto puro com medida de Lévy
explícita, e a parte aritmética consiste em saltos discretos induzidos por
dilatações e um termo de ``killing'' indexado por potências de primos até
\(\lambda^2\). O gerador auto-adjunto \(L_\lambda\) tem resolvente compacto e
estado fundamental estritamente positivo.

Este é o análogo contínuo do que o ES8-P14/16 faz discretamente. A tradução
para o projeto seria:

\begin{enumerate}
  \item \textbf{Identificar a matriz \(C_N\) como representação de Galerkin}
        da forma \(\E_\lambda\) numa base de funções de teste (por exemplo,
        exponenciais \(e_n(x)=\exp(i\pi n x/a)\) num intervalo de comprimento
        \(2a\)).
  \item \textbf{Provar que a forma discreta converge para a forma contínua}
        no sentido da \emph{convergência em forma} (norma
        \(\|\cdot\|_q^2 = \|\cdot\|^2 + q(\cdot)\)), não apenas entrada por
        entrada.
  \item \textbf{Provar que o limite é fechado e semibounded below}. Isso
        substitui as hipóteses ad hoc \texttt{SchurBounded} e
        \texttt{UniformFinitePositivity} do esqueleto Lean por uma propriedade
        estrutural da forma.
\end{enumerate}

\section{Onde o ES8-P16 já aponta o gargalo (e como contorná-lo)}

O ES8-P16 formaliza três obstáculos que precisam ser resolvidos antes de
qualquer extensão.

\subsection*{(a) Pivôs LDL não dão gap espectral uniforme}

O contraexemplo
\[
A_t=U_t^\top U_t,\qquad
U_t=\begin{pmatrix}1&t\\0&1\end{pmatrix},
\qquad
A_t=\begin{pmatrix}1&t\\t&t^2+1\end{pmatrix},
\]
tem pivôs \(1,1\) mas quociente de Rayleigh \(1/(t^2+1)\) no vetor \((-t,1)\).
A saída é \textbf{não usar LDL} como ferramenta de positividade em dimensão
infinita. Em vez disso, trabalhar diretamente com a forma fechada e usar o
\textbf{teorema KLMN} ou a \textbf{convergência em resolvente forte} de
compressões \(P_N M_\infty P_N\).

\subsection*{(b) O orçamento de erro por linha não escala}

A soma \(\sum_j E^{(N)}_{ij}\) já está no teto \(1/100\) para \(N=10\)
(linha 7: \(8333/10^6\)). Para \(N\to\infty\), é preciso substituir o orçamento
por linha por uma \textbf{norma de forma}:
\[
\|E_\infty\|_{\Lin(\D,\D^*)}\le \varepsilon,
\]
o que exige que \(E_\infty\) seja, por exemplo, Hilbert--Schmidt
(\(\sum_{i,j}E_{ij}^2<\infty\)). Os dados finitos (somas de Frobenius
\(\approx 1{,}2\times 10^{-11}\) em \(8\times 8\)) são compatíveis com isso,
mas \textbf{não provam} a convergência da cauda.

\subsection*{(c) A projeção \(B\) amplifica o erro pelo quadrado da norma}

Se \(y=Bx\) e \(|y^\top E y|\le \delta\|y\|^2\), então
\[
|x^\top B^\top E B x|\le \delta\,\|B\|_{\op}^2\,\|x\|^2.
\]
O ES8-P16 formaliza isso com o exemplo \(B=[2]\): erro bruto \(1/100\) vira
\(4/100\). Qualquer extensão precisa controlar
\(\sup_N\|B_N\|_{\op}^2\) ou, melhor, provar que a forma já está definida no
espaço ``físico'' sem projeção artificial.

\section{Rota de maior probabilidade: janelas compactas + extensão por convergência}

A literatura sobre \textbf{positividade de Weil em janelas compactas} é a mais
próxima do que o projeto pode atacar com rigor computacional:

\begin{itemize}
  \item Para suporte em
        \([-\tfrac12\log 2,\tfrac12\log 2]\), a positividade é clássica
        (Yoshida; Connes--Consani) e foi recentemente certificada com cotas
        explícitas via aritmética intervalar.
  \item Para suportes maiores (\([-0{,}8,0{,}8]\), autocorrelação \(1{,}6\)),
        há certificados de que \(Q(f)\ge 8{,}9\times 10^{-18}\|f\|^2\) para
        \textbf{todas} as \(f\) complexas de suporte \(1{,}6\).
  \item O perfil \(\lambda^*(L)=\inf Q(f)/\|f\|^2\) decai segundo uma
        \textbf{lei de Landau--Widom} duplamente exponencial: qualquer
        certificado ponto-a-ponto precisa resolver frequências até
        \(2\pi e^{A_L}\) com \(A_L\sim 4e^L\).
\end{itemize}

Isso significa que a extensão a \(L\to\infty\) \textbf{não pode ser feita por
refinamento de malha ou por aumento de \(N\)}. A única rota viável é:

\begin{enumerate}
  \item \textbf{Provar positividade numa janela fixa} (por exemplo,
        \(L=0{,}8\) ou maior) com certificado exato --- isso já está ao alcance
        do ferramental ES8-P16 (frações exatas \(+\) Lean).
  \item \textbf{Provar um princípio de propagação}: se a forma é positiva numa
        janela \([-L,L]\) e o erro de truncamento da cauda (primos e
        arquimediano) é controlado por uma cota que decai com \(L\), então a
        positividade se estende. A dificuldade é que a cota disponível
        \textbf{não decai} na taxa necessária --- ela é dominada pela cauda
        de zeros, que é justamente o que \RH{} controlaria.
  \item \textbf{Alternativa via operador}: construir um operador auto-adjunto
        \(L_\lambda\) cujo espectro seja real e positivo, e provar que seus
        autovalores coincidem com as ordenadas dos zeros. O resultado de forma
        de Dirichlet dá um candidato incondicional, mas a identificação
        espectral completa com os zeros de \(\zeta\) ainda não está provada.
\end{enumerate}

\section{O que o projeto pode fazer concretamente}

A Tabela~\ref{tab:roteiro} organiza o roteiro em ordem de custo/benefício.

\begin{table}[h!]
\centering
\small
\begin{tabular}{p{0.5cm} p{4.5cm} p{7cm}}
\toprule
\# & O que fazer & O que isso substitui no ES8-P16 \\
\midrule
1 & Formalizar nestedness e a unicidade da extensão
  & Já está no ES8-P16; consolidar como base.
    \texttt{extensionA}/\texttt{extensionB} mostram que o prefixo não basta
    --- é preciso uma \textbf{regra global} de definição. \\
2 & Substituir LDL por forma fechada
  & Provar que a forma discreta \(Q^{(N)}\) é restrição de uma forma contínua
    \(\E\) num espaço de Sobolev \(H^{\log}(-a,a)\) (ou similar).
    Elimina a dependência de pivôs e o contraexemplo \(A_t\). \\
3 & Provar fechamento e semiboundedness
  & Usar representação de Lévy--Khintchine para a parte arquimediana e soma
    finita de saltos para primos (até \(\lambda^2\)).
    Substitui \texttt{UniformFinitePositivity} por uma propriedade
    estrutural. \\
4 & Controlar a projeção \(B_N\)
  & Calcular ou limitar \(\sup_N\|B_N\|_{\op}^2\); se não for limitado,
    reformular a forma sem projeção.
    Elimina o fator \(\kappa\) não controlado do ES8-P16. \\
5 & Certificar positividade em janela fixa
  & Usar aritmética intervalar exata (já disponível no projeto) para
    reproduzir ou superar \(Q(f)\ge 8{,}9\times 10^{-18}\|f\|^2\) em
    \([-0{,}8,0{,}8]\). Fornece um ponto de ancoragem concreto para a
    extensão. \\
6 & Atacar a extensão \(L\to\infty\)
  & Provar que a cauda (primos com \(\log p>2L\) e arquimediano) é
    controlada por uma cota que não destrua a positividade. Aqui está o
    gargalo real. Nenhuma ferramenta atual do projeto cobre isso; exige
    análise fina da fórmula explícita. \\
\bottomrule
\end{tabular}
\caption{Roteiro técnico em ordem de custo/benefício.}
\label{tab:roteiro}
\end{table}

\section{Obstáculo fundamental que nenhuma formalização Lean resolve}

A lei de Landau--Widom mostra que a margem espectral \(\lambda^*(L)\)
\textbf{colapsa duplamente exponencialmente} com \(L\). Isso significa que,
para janelas grandes, a positividade se torna \textbf{numericamente
indetectável} por qualquer certificado ponto-a-ponto. Uma prova de \RH{} via
positividade de Weil exigiria um argumento \textbf{estrutural} que não passe
por resolução de frequências --- por exemplo, uma identificação espectral
exata entre um operador auto-adjunto e os zeros de \(\zeta\), ou uma prova de
que a forma é uma soma de quadrados por construção.

O ES8-P16 já registra corretamente que ``pivôs positivos não implicam gap
uniforme''. O próximo passo conceitual é reconhecer que \textbf{a
positividade de Weil não é uma propriedade perturbativa}: ela não sobrevive a
truncamentos arbitrários, e a extensão infinita não é obtida por ``mais do
mesmo''. É preciso:
\begin{enumerate}[label=(\roman*)]
  \item um operador global cuja positividade seja automática; ou
  \item um princípio de propagação não perturbativo que conecte janelas
        finitas ao todo.
\end{enumerate}

\section*{Resumo}

O caminho técnico existe e está mapeado na literatura (formas de Dirichlet,
janelas compactas certificadas, operadores de Hilbert--Pólya), mas ele
\textbf{não é uma continuação direta do ES8-P14/16}. É uma mudança de
paradigma: de ``verificar matrizes racionais finitas'' para ``provar
fechamento e positividade de uma forma em espaço de Hilbert''. O ES8-P16 é
útil precisamente por mostrar, com precisão cirúrgica, por que a ponte atual
não basta --- e por isolar, de forma verificável, os três obstáculos que
qualquer extensão precisa resolver.

\end{document}